\chapter{Máquina de Estados do Jogo}

\section{O \texttt{game\_state} global}

Uma única variável \texttt{game\_state} controla o comportamento do programa,
assumindo nove valores discretos. O \texttt{game\_tick} despacha lógica distinta para
cada estado; transições são disparadas por entrada do usuário, eventos da simulação
ou \texttt{setTimeout}s explícitos.

O \textbf{State Machine Diagram}, também conhecido como \emph{statechart}, comunica
simultaneamente os estados, transições e condições de guarda.

\diagfig{diagrams/state-machine.png}{Diagrama de Estados do \texttt{game\_state}}{state-machine}

\section{Catálogo de estados}

\begin{description}

  \item[\texttt{menu}] Tela de título. Exibe três opções (\textsc{JOGAR},
    \textsc{CRÉDITOS}, \textsc{SAIR}) navegáveis por setas ou WS, com confirmação por
    Enter ou Espaço. É o estado de entrada após o \emph{boot} e ponto de retorno
    garantido a partir de qualquer estado terminal.

  \item[\texttt{credits}] Tela de créditos com a equipe e referências de \emph{stack}.
    Implementa um \emph{grace period} de 300\,ms para ignorar o \emph{keyup} da
    própria tecla que abriu a tela. Qualquer tecla ou clique após esse intervalo
    retorna a \texttt{menu}.

  \item[\texttt{quit}] Tela de despedida ``SISTEMA DESLIGADO''. Tenta
    \texttt{window.close()}, que funciona no contexto Electron e falha silenciosamente
    no \emph{browser}. O \emph{fallback} para \texttt{menu} acontece via tecla ou
    clique.

  \item[\texttt{loading}] Estado transitório enquanto \texttt{regenerate\_atlas}
    aguarda o \texttt{Image} carregar. Acontece apenas no primeiro andar de cada
    sessão, pois chamadas subsequentes são síncronas com a imagem em cache.

  \item[\texttt{playing}] Estado central. O \texttt{game\_tick} executa as três fases
    do \emph{frame} (\texttt{update} $\rightarrow$ \texttt{collision} $\rightarrow$
    \texttt{render}), atualiza a HUD e checa o \emph{lock-and-clear}.

  \item[\texttt{paused}] Pausa por \textsc{ESC}. Sobrescreve os \emph{handlers} de
    teclado para ignorar tudo exceto novo \textsc{ESC}. O \texttt{game\_tick} faz
    \emph{early return} sem processar entidades nem renderizar; o \texttt{canvas}
    mantém o último \emph{frame} desenhado, e o \texttt{hud\_ctx} recebe uma camada
    semi-transparente com o texto ``PAUSADO''.

  \item[\texttt{shop}] Loja entre andares, ativada quando o jogador toca a escada em
    um andar não-final. Apresenta três \emph{upgrades} sorteados de um \emph{pool}
    persistente de cinco. \emph{Upgrades} comprados são removidos do \emph{pool} para
    o resto da \emph{run}.

  \item[\texttt{gameover}] Tela de derrota. Disparada pelo \texttt{setTimeout} de
    1500\,ms agendado em \texttt{entity\_player\_t.\_kill}. Suporta dois caminhos:
    \textsc{Enter} reinicia a \emph{run}, e \textsc{Esc} retorna ao menu.

  \item[\texttt{victory}] Tela de vitória. Disparada pelo \texttt{setTimeout} de
    2000\,ms em \texttt{entity\_boss\_t.\_kill}. Mostra estatísticas finais (tempo,
    kills, score, arma) e retorna ao menu por qualquer entrada.

\end{description}

\section{Transições críticas}

\subsection{\emph{Setup} de \emph{handlers} por estado}

Cada estado instala seu próprio conjunto de \emph{handlers}: \texttt{onkeydown},
\texttt{onkeyup}, \texttt{onmousedown} e \texttt{onmouseup}. A função
\texttt{bind\_input} restaura os \emph{handlers} de \emph{gameplay} ao entrar em
\texttt{playing}. Outros estados usam o \texttt{\_bind\_dismiss}, um \emph{helper}
que adiciona o \emph{grace period} e responde em \texttt{keyup} ou \texttt{mouseup},
ou implementam \emph{handlers} ad-hoc.

\subsection{Race conditions com \texttt{setTimeout}}

A morte do jogador agenda \texttt{show\_gameover\_screen} em 1500\,ms, e a morte do
\emph{boss} agenda \texttt{game\_victory} em 2000\,ms. Se as duas mortes acontecerem
no mesmo \emph{frame} (cenário plausível, em que a explosão final do \emph{boss}
mata o jogador), os dois \texttt{setTimeout}s ficam pendentes e disparam em janelas
diferentes. As transições estão protegidas por guardas idempotentes:

\begin{itemize}
  \item \texttt{show\_gameover\_screen} ignora se \texttt{game\_state} já é
    \texttt{'gameover'} ou \texttt{'victory'}.
  \item \texttt{game\_victory} ignora se \texttt{game\_state} já é \texttt{'victory'}.
\end{itemize}

A vitória, porém, \emph{sobrescreve} a derrota. A interpretação de design é direta:
se o jogador matou o \emph{boss}, o objetivo foi cumprido, independentemente de uma
morte contemporânea.

\subsection{Reset de \texttt{keys[]}}

Toda transição que entra em \texttt{playing} (a partir de \texttt{loading},
\texttt{shop} ou \texttt{gameover}) chama
\texttt{for (var k in keys) keys[k] = 0;} antes de devolver o controle. A linha evita
o problema de \emph{stuck keys}, em que uma tecla pressionada durante o estado
anterior ficaria registrada como ativa quando o jogo retomasse.
